\documentclass[11pt]{article}
\usepackage[utf8]{inputenc}
\usepackage{amsmath,amssymb}
\usepackage{graphicx}
\usepackage{booktabs}
\usepackage{authblk}
\usepackage[margin=1in]{geometry}
\usepackage[hidelinks]{hyperref}
\usepackage{setspace}

\title{Local update rules find stuck points, not minima:\\
\large free-boundary partitions, and a remedy}

\author[1]{Claes Johnson}
\affil[1]{KTH Royal Institute of Technology, SE-100 44 Stockholm, Sweden \\
\texttt{claesjohnson@gmail.com}}
\date{\today}

\begin{document}
\maketitle

\begin{abstract}
\noindent A class of solvers partitions space into domains, assigns a field to each, and evolves the
interfaces by a \emph{local} rule: a cell changes ownership when a neighbouring domain's field
exceeds its own, or when a locally computed velocity points that way. Level-set schemes with local
reinitialisation, Voronoi-partitioned solvers with density-balance interfaces, and free-boundary
formulations of multi-domain energy minimisation all have this structure.

Such a rule cannot minimise. It descends to the nearest configuration that no single admissible
change improves --- a \emph{stuck point} --- and which one it reaches depends on the initial
condition. This is reported here as a measured effect rather than a possibility. In a
multi-domain Coulomb energy minimisation, four independent geometries give converged energies that
differ by $0.005$--$0.05$ between seeds, between initial partitions, and between binary builds
differing only in floating-point rounding, against target quantities of $0.001$--$0.02$ in the same
units. The ambiguity exceeds the quantity sought by one to two orders in every case. Each run
converges cleanly; they converge to different places.

Four remedies were tried and none works: an interfacial energy penalty regularising the boundary, a
five-fold reduction of the interface time step, freezing the interface entirely, and changing to a
physical system without the first one's pathologies. The reason they fail is that this is a defect
of the \emph{search} and not of the energy functional, so nothing expressed in the functional
reaches it.

The remedy that does work is to stop searching. Interfaces in these problems are surfaces and can be
placed: specify the partition, hold it fixed, relax only the fields, and evaluate what a stated
configuration costs rather than asking where a boundary goes. Repeated runs then agree to every
digit recorded, and quantities that were previously irreproducible become well defined --- a
localised defect energy independent of system size to $5\%$, and an energy of delocalisation
independent of it likewise. The cost is that the family of admissible partitions must be
constructed rather than discovered, and three constructions are given, together with the invariants
that rejected two of them.
\end{abstract}

\section{The structure at issue}

Consider a domain of space divided into subdomains $\Omega_1,\dots,\Omega_N$ with
$\bigcup\Omega_i=\Omega$ and $\Omega_i\cap\Omega_j=\emptyset$, a field $u_i$ supported on each, and
an energy $E[\{u_i\},\{\Omega_i\}]$ to be minimised over both the fields and the partition. The
fields are handled by any standard relaxation. The partition is the difficulty, and the usual
treatment is an interface rule of the form
\begin{equation}\label{eq:rule}
\text{cell } c \text{ moves from } \Omega_i \text{ to } \Omega_j
\quad\text{when}\quad w_{ij}(c) < 0 ,
\end{equation}
where $w_{ij}$ is computed from the fields in a neighbourhood of $c$ --- a density balance, a local
energy difference, a level-set velocity. The rule is attractive because it is local, cheap, and
requires no global bookkeeping.

It is also, as an optimiser, a descent method whose moves are single-cell transfers. Its fixed
points are the configurations from which no single such transfer decreases the energy. These are
not the minima of $E$ over partitions, and there is in general no relation between the two sets
beyond inclusion of the latter in the former. That much is elementary. What is not obvious in
advance is the \emph{size} of the resulting ambiguity in a problem of practical interest, and
whether it can be controlled by the means one would ordinarily reach for.

\section{The measurement}

The setting is a minimisation of the Coulomb energy of $N$ non-overlapping unit charge densities,
one per domain, with the interfaces evolved by \eqref{eq:rule} using a density balance
$w=(\rho_i-\rho_j)/(\rho_i+\rho_j)$. Energies are in Hartree. Four geometries were run, each many
times, varying only the initial condition or the compiled binary:

\begin{center}
\small
\begin{tabular}{@{}llll@{}}
\toprule
geometry & what was varied & spread & target quantity\\
\midrule
ring of six sites      & the binary, physics identical & $0.005$ & $0.012$--$0.022$\\
chain of nine sites    & --- (never settles)           & $0.05$  & $\sim0.01$\\
two-centre transfer    & four initial partitions       & $0.031$ & $\sim0.001$\\
chain of five sites    & two initial partitions        & $0.013$ & $\sim0.001$\\
\bottomrule
\end{tabular}
\end{center}

\noindent Three remarks on the entries. The ring's spread appeared between two builds of the
\emph{same source} differing in an unrelated addition; recompilation perturbs floating-point
rounding, cells sitting near the threshold in \eqref{eq:rule} then fall the other way, and a
different stuck point follows. The chain of nine did not converge at all, oscillating between two
partitions indefinitely, which is the same phenomenon with a two-cycle instead of a fixed point. The
two-centre transfer was monotone in the initial displacement of one domain --- four seeds, four
energies, ordered --- which is the signature of a landscape of stuck points rather than of noise.

Each individual run is deterministic and reproducible: repeating one with the same seed and the same
binary returns the same energy to every digit. The irreproducibility is entirely across initial
conditions, and it exceeds the quantities being computed by one to two orders in all four cases.

\section{What does not fix it}

\paragraph{An interfacial penalty.} Adding $\tfrac{\beta}{2}\oint u^2\,dS$ to the energy gives the
interface a restoring force and is the natural regularisation. It behaves correctly as an energy: at
fixed initial condition it raises the total, linearly in $\beta$ at small $\beta$. It does not
reduce the spread. In the two-centre geometry the spread across seeds went from $0.031$ at
$\beta=0$ to $0.040$ at $\beta=0.05$ --- slightly worse. Spreading is an energy question and the
penalty answers it; which stuck point is reached is a search question, on which the penalty is
silent.

\paragraph{A smaller interface step.} Reducing the interface time step five-fold changes nothing.
Oscillation between two partitions looks like an overshoot and is not: the two configurations are
both fixed points of the fields, and the cycle is between them, not around one.

\paragraph{Freezing the interface.} Holding the partition fixed at its initial value removes the
search but also removes the answer, since the initial partition is then simply asserted. Applied to
the oscillating chain it left a residual variation of the same order, which localised the remaining
difficulty in the \emph{field} relaxation for that geometry --- a separate defect, and evidence that
the interface rule is not the only source of trouble in multi-domain solvers.

\paragraph{A different physical system.} The first geometries involved a weakly bound outer domain
whose field could lower its energy by spreading without limit, and one might attribute the spread to
that pathology. Replacing it with a system where every domain sits in a confining well removed the
pathology --- energies became physically sensible, with no runaway --- and the seed spread remained,
at $0.013$. The effect is not a property of the particular system.

\paragraph{And geometry cannot bound it either.} Confining domains to prescribed regions to prevent
spreading fails structurally when the partition must cover the domain: excluding a region from every
subdomain leaves it owned by whichever label the implementation defaults to. In our case that
produced a constrained minimum \emph{below} the unconstrained one --- impossible, and a useful
inequality to test against, since it exposed the error immediately.

\section{The remedy: specify the partition}

Interfaces in these problems are surfaces, and a surface can be placed. For a chain geometry the
interfaces are planes transverse to the axis, so a partition is fixed by the plane positions; more
generally any parametrised family will do. The procedure is:

\begin{enumerate}\itemsep2pt
\item construct a family of admissible partitions $\Omega(\lambda)$ parametrised by a coordinate
$\lambda$ carrying the physics of interest;
\item for each $\lambda$, fix the partition and relax the fields alone to convergence;
\item read $E(\lambda)$, and take from it whatever the problem requires --- a minimum, a barrier as
the amplitude over a period, a formation energy as a difference against a reference partition.
\end{enumerate}

Nothing is searched for, so nothing can be stuck. The question changes from ``where does the
boundary go'' to ``what does this configuration cost'', which is the question a barrier or a
formation energy actually poses. Repeated runs agree to every digit recorded --- verified three
times, at two system sizes, against the $0.005$--$0.05$ spreads above.

\paragraph{What it then delivers.} In the same multi-domain problem: a localised interface defect
costing $0.0064$ and $0.00675$ at two system sizes, i.e.\ independent of size to $5\%$ while the
total energy changes by $0.26$ between them --- a factor of $750$ in sensitivity; and an energy of
delocalisation of $0.111$ and $0.113$ at those two sizes, agreeing to $2\%$. Neither quantity was
obtainable at all by the free-boundary route, both being smaller than its ambiguity.

\paragraph{What it costs.} The family $\Omega(\lambda)$ must be constructed, and constructing it
correctly is not automatic. Three families were built for a single physical question and two were
wrong. One left a subdomain stranded against the boundary of $\Omega$ with no source in it; one
collapsed a subdomain to zero measure; the third was correct. Each was rejected before being run, by
checking an invariant the family had to satisfy --- that the ends of a chain be identical between
the configurations compared, that the subdomain count be preserved, that a translation by one period
return the original energy. \emph{Constructing the invariant is the substantive work}, and it is
worth doing before the computation rather than after, since a family that violates one produces
numbers that look entirely reasonable.

\section{Scope}

The finding applies to any scheme whose partition is advanced by a rule of the form
\eqref{eq:rule}, which includes local level-set reinitialisation, Voronoi- or power-diagram solvers
with locally computed interface velocities, and multi-domain energy minimisations of the kind
above. It does not apply to schemes that optimise the partition globally, nor to those in which the
partition is determined by the physics rather than sought --- a symmetry-fixed interface, for
example, is exact and costs nothing.

The remedy is available whenever the partitions of interest form a family that can be written down.
Where the partition itself is the unknown and no such family is available, the difficulty is real
and this note offers no solution to it; the honest course there is to report the ambiguity rather
than a converged energy, and to test for it, which requires only running the same problem twice from
different initial conditions. That test is cheap, it is not standard practice, and in four
geometries out of four here it would have been decisive.

\section{Conclusion}

A local interface rule is a descent method over single-cell moves, and its fixed points are stuck
points rather than minima. In a multi-domain Coulomb minimisation the resulting ambiguity was
$0.005$--$0.05$ Hartree across four geometries, one to two orders above the quantities being
computed, and was not removed by an interfacial penalty, a smaller step, a frozen interface, or a
change of physical system --- because it is a property of the search and not of the energy.
Specifying the partition and relaxing only the fields restores reproducibility to every recorded
digit and makes quantities computable that were not computable otherwise. The price is that the
family of partitions must be constructed and checked against an invariant, and the check must
precede the computation.

\end{document}
